\documentclass[../../../../main.tex]{subfiles}

\begin{document}

\section*{第1問}

［解］ $n = 0, 1, \dots, m-1, 1 \le m$ に対し，$f(m,n) = 5m + n$ と表される整数をもとめる。まず，

$25$ 以上の整数は全て $f(m,n)$ の形で示されることを示す。$m \ge 5$ とする。$m$ を固定すると，
$f(m,n)$ は $5m, 5m+1, \dots, 6m-1$ の全ての自然数をとる。
\[
\begin{cases}
f(m, m-1) = 6m - 1 \\
f(m+1, 0) = 5m + 5
\end{cases}
\]
において，$f(m, m-1) \ge f(m+1, 0) - 1 \iff m \ge 5$ だから，$m \ge 5$ ならば \dots を順に書き下していけば $25$ 以上の整数を全て表す。よって示された。

したがって $24$ 以下の自然数をしらべれば良い。$m \ge 5$ とすると $f(m,n) \ge 25$ だから，$m \le 4$ として書き下せば十分。
\[
\begin{cases}
m = 1 & \dots \ 5 \\
m = 2 & \dots \ 10, 11 \\
m = 3 & \dots \ 15, 16, 17 \\
m = 4 & \dots \ 20, 21, 22, 23
\end{cases}
\]
だから，これに含まれない $24$ 以下の自然数を列挙して
\[
1, 2, 3, 4, 6, 7, 8, 9, 12, 13, 14, 18, 19, 24
\]

\newpage

\section*{第2問}

［解］ 長針，短針の長さを $R, r$ とする。

\begin{center}
\begin{tikzpicture}[scale=0.8, >=stealth]
  % 2時
  \begin{scope}[shift={(0,0)}]
    \draw[thick] (0,0) circle (1);
    \draw[thick] (0,0) circle (0.5);
    \draw[->, thick] (0,0) -- (0, 0.9);
    \draw[->, thick] (0,0) -- (60:0.45);
    \node at (0, -1.4) {2時};
    \node at (75:1.2) {$\frac{\pi}{3}$};
  \end{scope}
  % 2時半
  \begin{scope}[shift={(3.5,0)}]
    \draw[thick] (0,0) circle (1);
    \draw[thick] (0,0) circle (0.5);
    \draw[->, thick] (0,0) -- (0, -0.9);
    \draw[->, thick] (0,0) -- (75:0.45);
    \node at (0, -1.4) {2時半};
    \node at (105:1.2) {$\frac{7}{12}\pi$};
  \end{scope}
  % 4時
  \begin{scope}[shift={(7,0)}]
    \draw[thick] (0,0) circle (1);
    \draw[thick] (0,0) circle (0.5);
    \draw[->, thick] (0,0) -- (0, 0.9);
    \draw[->, thick] (0,0) -- (-30:0.45);
    \node at (0, -1.4) {4時};
    \node at (30:1.2) {$\frac{2}{3}\pi$};
  \end{scope}
\end{tikzpicture}
\end{center}

題意及び余弦定理から，
\[
\begin{cases}
4^2 = R^2 + r^2 - 2Rr \cos \frac{\pi}{3} = R^2 + r^2 - Rr & \dots \text{①} \\
6^2 = R^2 + r^2 - 2Rr \cos \frac{7}{12}\pi &
\end{cases}
\]

又，もとめる長さ $l\ (>0)$ として
\[
l^2 = R^2 + r^2 - 2Rr \cos \frac{2}{3}\pi = R^2 + r^2 + Rr \quad \dots \text{②}
\]

ここで，$\cos \frac{7}{12}\pi < 0$ から，
\[
\cos \frac{7}{12}\pi = -\sqrt{\frac{1 + \cos \frac{7}{6}\pi}{2}} = -\sqrt{\frac{1 - \frac{\sqrt{3}}{2}}{2}} = -\frac{\sqrt{3}-1}{2\sqrt{2}}
\]
だから，①より，
\begin{align*}
36 &= R^2 + r^2 + 2Rr \frac{\sqrt{3}-1}{2\sqrt{2}} \\
&= R^2 + r^2 + \frac{\sqrt{3}-1}{\sqrt{2}} Rr \quad \dots \text{③}
\end{align*}

①，③の辺々引いて，
\[
20 = \left( \frac{\sqrt{6}-\sqrt{2}}{2} + 1 \right) Rr
\]
\[
\therefore Rr = \frac{40}{\sqrt{6}-\sqrt{2}+2} \quad \dots \text{④}
\]

②に代入して
\begin{align*}
l^2 &= (16 + Rr) + Rr \quad (\because \text{①}) \\
&= 16 + 2Rr \\
&= 16 + \frac{80}{\sqrt{6}-\sqrt{2}+2} \quad (\because \text{④}) \\
&= 16 + 20(\sqrt{3}+1-\sqrt{2}) \\
&= 20(\sqrt{3}-\sqrt{2}) + 36 \quad \dots \text{⑤}
\end{align*}

ここで，$(6.5)^2 < l^2 < (6.6)^2$ を示す。
\begin{align*}
(1.73)^2 < 3 < (1.733)^2 \quad &\therefore 1.73 < \sqrt{3} < 1.733 \\
(1.41)^2 < 2 < (1.42)^2 \quad &\therefore 1.41 < \sqrt{2} < 1.42
\end{align*}
だから，⑤に代入して
\[
20(1.733 - 1.42) + 36 < l^2 < 20(1.74 - 1.41) + 36
\]
\[
\therefore 42.26 < l^2 < 42.6
\]
一方，$(6.5)^2 = 42.25$, $(6.6)^2 = 43.56$ だから，⑤より，
\[
(6.5)^2 < l^2 < (6.6)^2
\]
よって示された。$l > 0$ より，
\[
6.5 < l < 6.6
\]
つまり，
\[
l \fallingdotseq 6.5
\]
である。

\newpage

\section*{第3問}

［解］ $C: (x - 2a)^2 + (y - a)^2 = 5a^2 - 20a + 25$ とおく。

$(\text{与式}) \iff (-4x - 2y + 20)a + (x^2 + y^2 - 25) = 0$ だから，これが $a$ についての恒等式の時，つまり
\[
\begin{cases}
-4x - 2y + 20 = 0 \\
x^2 + y^2 = 25
\end{cases}
\iff (x,y) = (3,4), (5,0)
\]
の時，与式は $a$ によらず成り立つから，もとめる $2$ 点はこれである。又，$2$ 円の中心間距離 $d$ は $d = \sqrt{(2a)^2 + a^2} = \sqrt{5}|a|$ である。

\paragraph{1. 外接する時}
\[
\sqrt{5} + \sqrt{5}\sqrt{a^2 - 4a + 5} = \sqrt{5}|a|
\]
\[
\sqrt{a^2 - 4a + 5} = |a| - 1
\]
$|a| \ge 1$ のもとで $2$ 乗して，整理すると $a = 2$ を得る。

\paragraph{2. 内接する時}
\[
| \sqrt{5} - \sqrt{5}\sqrt{a^2 - 4a + 5} | = \sqrt{5}|a|
\]
両辺 $0$ 以上から $2$ 乗して
\[
a^2 - 4a + 5 + 1 - 2\sqrt{a^2 - 4a + 5} = a^2
\]
\[
\sqrt{a^2 - 4a + 5} = 3 - 2a
\]
$a \le \frac{3}{2}$ のもとで $2$ 乗して
\[
a^2 - 4a + 5 = 4a^2 - 12a + 9
\]
\[
3a^2 - 8a + 4 = 0
\]
\[
(3a - 2)(a - 2) = 0 \quad \dots \quad \therefore a = \frac{2}{3}, 2
\]
$a \le \frac{3}{2}$ をみたすのは $a = \frac{2}{3}$

以上から $a = \frac{2}{3}, 2$

\newpage

\section*{第4問}

［解］ 与式を $f(x,y)$ とおく。$S = \sin x, C = \cos x$ とおく。
\[
f(x,y) = 2S + \begin{pmatrix} 1 \\ \sqrt{3} \end{pmatrix} \cdot \begin{pmatrix} \cos y \\ \sin y \end{pmatrix} \cdot C \quad \dots \text{①}
\]
コーシー・シュワルツから
\[
-2 \le \begin{pmatrix} 1 \\ \sqrt{3} \end{pmatrix} \cdot \begin{pmatrix} \cos y \\ \sin y \end{pmatrix} \le 2 \quad \left(\text{左側は } y = \frac{4}{3}\pi, \text{右は } y = \frac{\pi}{3}\right) \quad \dots \text{③}
\]
である。又，以下のように場合分けする。

ア：$0 \le x \le \pi/2$，イ：$\pi/2 \le x \le \pi$，ウ：$\pi \le x \le \frac{3}{2}\pi$，エ：$\frac{3}{2}\pi \le x \le 2\pi$

下表を得る。

\begin{center}
\begin{tabular}{|c|c|c|}
\hline
& ア & イ \\
\hline
$\max f(x,y)$ & $2(C+S) \le 2\sqrt{2} \quad (x=\pi/4 \text{ で等号})$ & $2(S-C) \le 2\sqrt{2} \quad (x=\frac{3}{4}\pi)$ \\
\hline
$\min f(x,y)$ & $2(S-C) \ge -2 \quad (x=0)$ & $2(S+C) \ge -2 \quad (x=\pi)$ \\
\hline
\end{tabular}

\medskip

\begin{tabular}{|c|c|c|}
\hline
& ウ & エ \\
\hline
$\max f$ & $2(S-C) \le 2 \quad (x=\pi)$ & $2(S+C) \le 2 \quad (x \to \pi)$ \\
\hline
$\min f$ & $2(C+S) \ge -2\sqrt{2} \quad (x=\frac{5}{4}\pi)$ & $2(S-C) \ge -2\sqrt{2} \quad (x=\frac{7}{4}\pi)$ \\
\hline
\end{tabular}
\end{center}

よって
\[
\begin{cases}
\max = 2\sqrt{2} & \left( (x,y) = \left(\frac{\pi}{4}, \frac{\pi}{3}\right), \left(\frac{3}{4}\pi, \frac{4}{3}\pi\right) \right) \\[1ex]
\min = -2\sqrt{2} & \left( (x,y) = \left(\frac{5}{4}\pi, \frac{\pi}{3}\right), \left(\frac{7}{4}\pi, \frac{4}{3}\pi\right) \right)
\end{cases}
\]

\newpage

\section*{第5問}

［解］ 漸化式から
\begin{align*}
f_1(x) &= \frac{1}{2h} \int_{x-h}^{x+h} \cos x \, dx = \frac{1}{2h} [\sin(x+h) - \sin(x-h)] \\
&= \frac{1}{2h} \cdot 2\sin h \cos x = \frac{\sin h}{h} \cos x = \frac{\sin h}{h} f_0(x)
\end{align*}
だから，これをくり返し用いて，$p = \frac{\sin h}{h}$ として
\[
f_n(x) = p^n \cos x
\]
だから
\begin{align*}
\sum_{k=1}^n f_k(x) &= p \cos x \frac{1 - p^n}{1 - p} \xrightarrow{n \to \infty} \frac{p}{1 - p} \cos x \quad (\because |p| < 1) \\
&= \frac{\sin h}{h - \sin h} \cos x
\end{align*}

\newpage

\section*{第6問}

［解］ 平均値の定理から，$h(x) = \frac{f(x) - f(a)}{x - a}$ に対し，$h(x) = f'(c)$ なる $c$ が $a < c < x$ にある。 \dots ①

又，$f''(x) > 0$ から $f'(x)$ は単調増加。 \dots ②

$a < x < b$ のとき，
\begin{align*}
h'(x) > 0 &\iff f'(x)(x-a) - [f(x) - f(a)] > 0 \\
&\iff f'(x) > \frac{f(x) - f(a)}{x - a} = f'(c) \quad \dots \text{③}
\end{align*}
だが，②及び $c < x$ から③は明らか。よって $h'(x) > 0$ つまり $h(x)$ は単調増加。

\end{document}
